%% =================================================================
%% Wei, Mei, Zhao, Xiao, Sun, Zhang, Guo.
%% "Nondestructively identifying the mechanical behavior of soft
%%  tissues using surface deformation with an explicit inverse
%%  approach."
%% Appl. Math. Modelling 134 (2024) 126--147.
%%
%% Reconstruction of page 5 (printed p. 130) as study notes.
%% Prose: extracted verbatim from the PDF text layer via pdftotext
%%        (line breaks and bracket-wrapping artifacts cleaned up).
%% Math:  hand-transcribed from the rendered page image — Eq. (15),
%%        the TVD (total variation diminishing) regularizer.
%% Figures 3 & 4: placeholder boxes with captions + visual descriptions.
%% =================================================================

\documentclass[11pt]{article}

\usepackage[margin=1in]{geometry}
\usepackage{amsmath, amssymb}
\usepackage{mathrsfs}
\usepackage{bm}
\usepackage{xcolor}
\usepackage{fancyhdr}

\pagestyle{fancy}
\fancyhf{}
\lhead{\small Y.~Mei et al.}
\rhead{\small Applied Mathematical Modelling 134 (2024) 126--147}
\cfoot{\small \thepage}
\renewcommand{\headrulewidth}{0.4pt}

\setcounter{page}{130}
\setcounter{equation}{14}  % next equation is (15)

\begin{document}

% ---------- Figure 3 placeholder ----------
\noindent
\begin{center}
\fbox{\parbox{0.92\textwidth}{\centering\vspace{1em}
\textbf{[Figure 3 — placeholder]}\\[0.8em]
Square light-teal domain with background shear modulus $\mu^{0}$,
containing an irregular pink inclusion region with shear modulus
$\mu^{1}$ and boundary $\partial\Omega^{1}$. Eight control points
$\mathbf{P}^{1}_{1}, \mathbf{P}^{1}_{2}, \ldots, \mathbf{P}^{1}_{8}$
are distributed around the boundary, connected to each other by
dashed blue segments (the control polygon) and smoothly interpolated
by an orange closed B-spline curve. Eight radial distances
$d^{1}_{1}, d^{1}_{2}, \ldots, d^{1}_{8}$ connect the inclusion
center $(x^{1}_{0}, y^{1}_{0})$ to each control point, illustrating
the polar parameterization used in Eq.~(7) and Eq.~(10).
\vspace{1em}}}
\end{center}
\vspace{0.3em}
\begin{center}
\textbf{Fig. 3.} The nonhomogeneous shear modulus distribution of an
inclusion embedded composite.
\end{center}

\vspace{0.8em}

% ---------- Figure 4 placeholder ----------
\noindent
\begin{center}
\fbox{\parbox{0.92\textwidth}{\centering\vspace{1em}
\textbf{[Figure 4 — placeholder]}\\[0.8em]
Graph of $H(z) = 0.5 \times (1 + \tanh(z/\varepsilon))$ (the smoothed
Heaviside from Eq.~(13)) plotted over $z \in [-0.5, 0.5]$, with
$H(z) \in [0, 1]$ on the vertical axis. Four curves show the effect
of decreasing $\varepsilon$:\\[0.3em]
\begin{tabular}{ll}
\textcolor{red}{$\varepsilon = 0.01$} & sharpest transition (nearly a true step) \\
\textcolor{green!60!black}{$\varepsilon = 0.05$} & moderate transition \\
\textcolor{blue}{$\varepsilon = 0.10$} & smoother transition \\
\textcolor{orange}{$\varepsilon = 0.15$} & smoothest transition
\end{tabular}\\[0.3em]
All four curves pass through $(0, 0.5)$. As $\varepsilon \to 0$ the
approximation converges to the true Heaviside; as $\varepsilon$ grows
the transition band widens.
\vspace{1em}}}
\end{center}
\vspace{0.3em}
\begin{center}
\textbf{Fig. 4.} An approximation of Heaviside function.
\end{center}

\vspace{1em}

% ---------- Eq. 15 ----------
\noindent
\begin{equation}
\text{Reg}(\mu) = \dfrac{1}{2}\,\alpha \int_{\Omega}
\sqrt{|\nabla\mu|^{2} + c^{2}}\, d\Omega
\end{equation}

\vspace{0.3em}

% ---------- Prose after Eq. 15 ----------
\noindent
For the layered structure,
$\mathbf{r}^i = \left[r^i_1(t_1),\, r^i_2(t_2),\, \cdots,\, r^i_n(t_n)\right]$
is the geometric parameter vector representing the distance between
the $i$-th interface and $C^0$ at the location $t_k$. For the
structure with inclusions embedded, the geometric parameter vector
$\mathbf{r}^i = \left[r^i_1,\, r^i_2,\, \cdots,\, r^i_n\right]
= \left[x^i_0,\, y^i_0,\, d^i_1,\, \cdots,\, d^i_{n-2}\right]$
is consisting of the coordinates of the center and distances between
the control points to the center of $i$-th component. In Eq.~(14),
the discrepancy between the computed displacement field and measured
displacement field on partial surface $\partial\Omega_{i_{\text{meas}}}$
is minimized in $L$-$2$ norm. The computed displacement field is
obtained by solving a forward problem using the finite element method
[32]. Due to the influence of noise in measurements, we introduce a
total variation diminishing regularization (TVD) term (the second
term in Eq.~(14)), and the specific form is given by Eq.~(15).
$\alpha$ is the regularization factor which controls the significance
of the regularization on the objective function. $c$ is a small
constant to avoid singularities when evaluating the gradient of the
regularization term with respect to optimization variables. It should
be noted that there exist alternative forms for the regularization
term apart from the domain integral forms. Nevertheless, we have
chosen to adhere to the use of domain integrals due to practical
considerations. In particular, the explicit inverse solver has been
developed based on the established framework of the existing
implicit inverse solver. This continuity in methodology ensures
consistency and facilitates seamless integration between the two
solvers [33]. The specific value of the regularization factor is
chosen based on the smoothness criteria [22].

\end{document}
